\section{信息论时间：柱集收缩与熵率箭头}\label{sec:info-time}

\subsection{历史一致集合与滤过（D6.1）}

对稳定输出前缀 $x_0^T:=(x_0,\dots,x_T)$，定义历史一致集合
\[
C_T(x_0^T):=\left\{\omega\in\{0,1\}^{T+m}:\ \forall 0\le t\le T,\ \Fold_m(\omega^{(m)}_t)=x_t\right\}.
\]
该集合刻画在给定输出约束下仍不可区分的原始读出候选。随 $T$ 增大，约束增加，集合在典型情形下收缩。

\subsection{时间箭头的协议化定义（D6.2）}

定义协议时间箭头为：
\begin{itemize}
  \item 若存在一族可审计的后验规模量（例如一致集合的诱导测度、条件熵、或最优编码长度）随 $T$ 单调收缩，则以 $T$ 增大方向定义时间箭头。
\end{itemize}
该定义完全位于协议层，不引入对象层时间坐标。

\begin{remark}[“未来”的对象化口径]
在本文中，“未来”严格对应条件分布的不确定性规模，例如
$
H(x_{t+1}\mid x_{\le t})
$
或更一般的后验柱集测度收缩指标；不将余项实体化为额外本体对象。
\end{remark}

\subsection{Shannon--McMillan--Breiman 机制（T6.3，H6.3）}

在遍历测度保持系统与有限生成分割条件下，Shannon--McMillan--Breiman（SMB）定理给出典型柱集测度的指数收缩率等于度量熵。本文使用如下模板：
\begin{itemize}
  \item \textbf{（H6.3）}稳定化读出 $(x_t)$ 由某遍历过程诱导，并且所用分割在相应因子系统上为生成分割；
  \item \textbf{（T6.3）}则典型 $x_0^T$ 的柱集测度满足
  \[
  -\frac{1}{T}\log \mu([x_0^T])\to h_\mu,
  \]
  其中 $h_\mu$ 为 Kolmogorov--Sinai 熵。
\end{itemize}
由此，“时间箭头”可由熵率控制的指数收缩机制给出严格化表达（附录 \ref{app:smb}）。

